% Source: https://artofproblemsolving.com/community/c3h1715915_junior_regional_olympiad__fbh_2015_grade_8_problem_3
% Stable identity: aops-topic-1715915-q01
Let $E'$ be the reflection of $E$ in $D$, so that $D$ is the midpoint of $EE'$.

\begin{asy}
usepackage("amsmath");
size(6cm);
defaultpen(fontsize(10pt));
pair A=(1.5,4),B=(0,0),C=(6,0),D=(B+C)/2;
pair Ep=B+0.5*(A-B);
pair F=extension(D,D+rotate(90)*(Ep-D),A,C);
pair Ep2=2*D-Ep;
void ra(pair P,pair U,pair V,real s){draw(P+s*unit(U-P)--P+s*unit(U-P)+s*unit(V-P)--P+s*unit(V-P),linewidth(0.5));}
draw(A--B--C--cycle,linewidth(0.9));
draw(Ep--D--F,linewidth(0.8));draw(Ep--F,linewidth(0.8));
draw(D--Ep2,blue+dashed);draw(C--Ep2,blue+dashed);draw(F--Ep2,blue+dashed);
ra(D,Ep,F,0.3);
dot(D);dot(Ep);dot(F);dot(Ep2,blue);
label("$A$",A,N);label("$B$",B,SW);label("$C$",C,E);label("$D$",D,dir(-100));
label("$E$",Ep,W);label("$F$",F,NE);label("$E'$",Ep2,SE,blue);
label("\mbox{}",(min(currentpicture,user=true)+max(currentpicture,user=true))/2,black); // keeps the text after the figure black when LaTeX typesets the labels inline
\end{asy}

\textbf{$CE'=BE$.} The segments $BC$ and $EE'$ bisect each other at $D$, so $EBE'C$ is a parallelogram (or, directly, triangles $DBE$ and $DCE'$ are congruent by SAS with the vertical angles at $D$). Hence $CE'=BE$.

\textbf{$FE'=FE$.} Since $\angle EDF=90^\circ$ and $D$ is the midpoint of $EE'$, the line $DF$ is the perpendicular bisector of $EE'$, so $FE'=FE$.

\textbf{Triangle inequality.} In triangle $FCE'$,
\[
CF+CE'\ge FE',\qquad\text{that is,}\qquad CF+BE\ge EF ,
\]
with equality only if $C$ lies on segment $FE'$. That would put $E'$ on line $AC$. But $CE'\parallel BE$, i.e.\ $CE'$ is parallel to $AB$, and the line through $C$ parallel to $AB$ meets line $AC$ only at $C$; so $E'=C$ and $E=B$, which is excluded. Therefore $BE+CF>EF$. $\blacksquare$
